\section{Introduction}

\subsection{Why Self-Healing for Vulnerability Remediation}

A self-healing vulnerability remediation framework detects vulnerable
hosts, decides which to patch, executes the patches, and rolls back
when the post-action state is unhealthy --- all without human
intervention beyond exception handling. Commercial implementations
exist (Microsoft Defender automated investigation; Tanium; IBM SOAR);
the academic literature has focused on the upstream components
(scoring, prioritization) and not on the integrated closed-loop
system with pre-registered safety bounds.

Papers 3--9 of this research sequence developed those upstream
components under a consistent pre-registration discipline:
HygieneBench (Paper~3)~\cite{paper3hygienebench} for the synthetic
telemetry substrate; HygienePrio (Paper~4)~\cite{paper4hygieneprio}
for the scoring rule;
Paper~7~\cite{paper7online} for the deployable online calibration
recipe; Paper~6~\cite{paper6capacity} for the capacity-arrival
operating regime; Paper~9 Theorem~1~\cite{paper9selftraj} for the
structural bound on closed-loop signal exhaustion; and the real
EPSS/KEV public corpus (released with the external-validity section
of Paper~4, \S X~\cite{paper4hygieneprio})
for realistic CVE attribute distributions.

This paper integrates those components into AutoHeal --- a
seven-stage closed-loop pipeline with explicit pre-registered safety
bounds --- and evaluates it on a frozen 2{,}700-row sweep.

\subsection{Pre-Registered Hypotheses}

The four hypotheses locked before evaluation are:

\textbf{H1 (Coverage):} AutoHeal remediates $\geq 80\%$ of high-EPSS
pairs by Window 6 at moderate capacity ($K \in \{50, 100\}$).
$K = 200$ is excluded by Paper~9 Corollary~4: closed-loop
calibration is structurally unable to escape signal exhaustion at
that capacity.

\textbf{H2 (Safety):} Per-window rollback rate $\leq 5\%$ at every
cell-seed-window triple. The hard-stop threshold (10\%) is twice
the H2 threshold; H2 is the operational target.

\textbf{H3 (MTTR Reduction):} AutoHeal's mean time to remediation
is $\leq 50\%$ of the Human-in-loop baseline's MTTR at moderate
capacity.

\textbf{H4 (Per-Pair Dominance):} AutoHeal beats Human-in-loop in
$\geq 80\%$ of (cell, seed, window) triples at $K \leq 100$. The
protocol registered this hypothesis on Precision@50; the analysis
code evaluated per-window coverage dominance. This deviation is
disclosed in \S\ref{sec:results}.

\subsection{Contributions}

\begin{itemize}
\item \textbf{C1 --- A pre-registered self-healing architecture.}
AutoHeal's seven-stage pipeline is documented in \S\ref{sec:architecture}
with triage thresholds, failure-mode distribution, and safety bounds
locked before evaluation. The architecture is the artifact;
re-implementations on different scoring rules are direct.

\item \textbf{C2 --- A 2{,}700-row frozen evaluation.}
\S\ref{sec:results} reports outcomes on the EEHDA fleet with real
CVE attributes for 25 evaluation seeds across 12 windows and three
capacities.

\item \textbf{C3 --- H1 supported: autonomous remediation works at
moderate capacity.} AutoHeal eventually remediates $97.8\%$ of
high-EPSS pairs by Window 6, exceeding the pre-registered $80\%$
threshold by ${\approx}18$~pp.

\item \textbf{C4 --- H2 rejected: the safety bound is exceeded
substantively.} The hard-stop bound is flagged 134 times across the
sweep (14.9\% of windows); cascading failures detected three
times. \emph{The detection mechanism works as designed} --- it
correctly flags windows where conditions warrant escalation
(in the simulated evaluation the flag suppresses new-CVE intake
for the affected window; enforcement by hard stop with fallback
was not exercised, see \S\ref{sec:architecture} and
\S\ref{sec:threats}). The H2 tolerance ($5\%$) was set assuming the pre-registered
failure-mode distribution would produce stable per-window
rollback rates; in practice small AUTO buckets at the
conservative triage threshold amplify per-window rates beyond the
H2 tolerance.

\item \textbf{C5 --- H4 rejected: structural-not-throughput parity
with human-in-loop.} Under the pre-registered conservative triage
thresholds, AutoHeal's per-window AUTO bucket averages
$\approx 27$ pairs at $K = 50$ (of which $\approx 19$ are acted
per window), comparable to Human-in-loop's
unconditional 30. AutoHeal trades throughput for safety-bounded
automation rather than the other way around.

\item \textbf{C6 --- H3 not analysable.} The MTTR instrumentation
records disclosure at remediation time rather than detection
time, producing degenerate cell-mean MTTR estimates ($0.0$ for
both AutoHeal and Human-in-loop). We report the issue honestly
per the protocol's stop rule and pre-register a corrected
instrumentation for the follow-up evaluation.
\end{itemize}

\subsection{Relationship to Prior Papers}

This paper is the synthesis paper of the VulnPrio sequence. Papers~1
and 3--9 produced components and bounds; AutoHeal integrates them
under a pre-registered architecture. The mixed outcome of the
four hypotheses --- H1 supported, H2 rejected, H4 rejected, H3 not
analysable --- demonstrates that the components fit together in a
working closed-loop system whose failure modes are predicted by
the prior papers' findings (Paper~9 Corollary~4 predicts the
K=200 exclusion; Paper~7 predicts the K=50/100 calibration recipe
choice).
